\documentclass{article}
\usepackage{amsmath}
\usepackage{amssymb}
\usepackage{amsthm}
\usepackage{geometry}
\usepackage{graphicx}
\usepackage{hyperref}
\usepackage{natbib}
\geometry{margin=1in}

\title{Mind at Work: Detecting Agentic Intelligence through the Effective Assembly Index and Praxicon Framework}
\author{Flyxion}
\date{September 10, 2025}

\begin{document}

\maketitle

\begin{abstract}
This essay presents a comprehensive framework for detecting mind-like processes by integrating the Effective Assembly Index for Mind Recognition (EAIMR) with the praxicon concept from cognitive science. EAIMR quantifies the improbability of a system’s structure arising from unguided processes, adjusted for efficiency and amplified by coherence, while the praxicon provides a structured repository of action representations bridging perceptual, semantic, and motor domains. By combining these, we propose a robust method to identify agentic intelligence in neural systems, collective organisms, ecosystems, and artificial intelligence, avoiding anthropocentric biases. The framework is formalized using information theory, graph theory, and category theory, with applications in neuroscience, AI alignment, ecological cognition, and cross-cultural studies. We contrast this approach with competing theories, provide detailed case studies, and discuss ethical implications for recognizing non-human minds.
\end{abstract}

\section{Introduction}
\label{sec:introduction}

Detecting mind-like processes across diverse systems—biological, ecological, and artificial—remains a critical challenge in cognitive science, artificial intelligence (AI), and philosophy of mind. The Effective Assembly Index for Mind Recognition (EAIMR) offers a quantitative measure of improbable, goal-directed assembly, distinguishing agentic systems from those driven by unguided physical processes \citep{walker2013}. Concurrently, the praxicon, a cognitive architecture for action representation, bridges perceptual, semantic, and motor domains, providing a structured framework for understanding embodied intelligence \citep{pastra2011}. This essay integrates EAIMR with the praxicon to create a unified approach for detecting mind-like agency, applicable to neural circuits, termite mounds, forest ecosystems, and AI systems. By grounding this integration in information theory, graph theory, and category theory, we address interdisciplinary questions in AI alignment, ecological cognition, and interspecies ethics, while avoiding narrativity biases that privilege human-like communication.

\section{Historical and Conceptual Foundations}
\label{sec:foundations}

\subsection{Assembly Theory and Complexity}
Assembly theory, pioneered by researchers like Hazen, Cronin, and Walker, quantifies the complexity of systems by measuring the minimum steps required to construct them from elemental primitives \citep{hazen2007, cronin2016, walker2013}. Originating in studies of biocomplexity and the origins of life, it emphasizes functional information and selection processes that produce improbable structures. EAIMR extends this by adjusting for reuse, parallelism, degeneracy, and coherence, making it a candidate for detecting mind-like processes \citep{adami2012}.

\subsection{Praxicon: Action Representation in Cognition}
The praxicon, introduced in neuropsychological models of limb praxis, is a long-term memory repository for familiar actions, comprising input (recognition) and output (production) components \citep{peigneux2004, mantovani2010}. Recent work extends it to computational models and robotics, grounding language in motor primitives \citep{pastra2011, panunzi2018}. It interfaces with semantic systems, enabling mappings between abstract concepts and embodied actions, with neural substrates in parietal-premotor circuits \citep{vannuscorps2023}.

\subsection{Competing Approaches}
EAIMR-praxicon contrasts with other mind-detection frameworks:
- \textbf{Integrated Information Theory (IIT)} \citep{tononi2004} quantifies consciousness via information integration but struggles with non-neural systems.
- \textbf{Free Energy Principle (FEP)} \citep{friston2010} models agency as minimization of prediction error, yet lacks specificity for action representation.
- \textbf{Enactive Cognition} \citep{barandiaran2009} emphasizes sensorimotor loops but does not quantify structural improbability as EAIMR does.

\section{Formal Definitions}
\label{sec:formal}

\subsection{EAIMR Definition}
For a system \(X\) over alphabet \(\Sigma\), EAIMR is defined as:

\[
\text{EAIMR}(X) = \frac{A(X)}{R(X) \cdot P(X) \cdot D(X)} \cdot C(X),
\]

where:
- \(A(X)\): Raw Assembly Index, minimum steps to build \(X\) \citep{walker2013}.
- \(R(X)\): Reuse factor, reflecting modular components.
- \(P(X)\): Parallelism factor, for concurrent assembly.
- \(D(X)\): Degeneracy factor, number of functionally equivalent configurations.
- \(C(X)\): Coherence factor, measuring goal-directed integration.

High EAIMR indicates low likelihood under null (mindless) dynamics \(P_0(X)\), suggesting agentic processes \citep{adami2012}.

\subsection{Praxicon Definition}
The praxicon is a bipartite graph \(G_{\mathcal{P}} = (V_A, V_S, E)\), where \(V_A\) are action nodes (motor engrams), \(V_S\) are semantic tags, and \(E\) are edges mapping actions to meanings \citep{pastra2011}. Input praxicon handles recognition, output praxicon handles production, interfacing via an action-semantic system \citep{peigneux2004}.

\section{Mathematical Integration}
\label{sec:integration}

\subsection{Information-Theoretic Foundations}
Let \(P_0(X)\) and \(P_1(X)\) be null and mind-like generative distributions. Self-information under null:

\[
I_0(X) = -\log P_0(X).
\]

Likelihood ratio:

\[
\Lambda(X) = \frac{P_1(X)}{P_0(X)} = \exp(\log P_1(X) - \log P_0(X)).
\]

EAIMR correlates with high \(\Lambda(X)\) and \(I_0(X)\). Praxicon reduces action entropy:

\[
\Delta H_{\mathcal{P}} = H(\mathcal{P}_0) - H_{\text{obs}}(\mathcal{P}),
\]

where \(\mathcal{P}_0\) is a null action distribution \citep{shannon1948}. Then,

\[
C(X) \approx 1 + \alpha \frac{\Delta H_{\mathcal{P}}}{H(\mathcal{P}_0)}.
\]

Kolmogorov complexity \(K(X)\) approximates \(A(X)\), with praxicon reducing \(K_{\text{plan}}(X)\) via stored engrams \citep{kolmogorov1965}.

\subsection{Graph-Theoretic Formalism}
Model praxicon as a bipartite graph \(G_{\mathcal{P}}\). EAIMR highlights improbable subgraph regularities (e.g., low-entropy action clusters). Define coherence via graph stability:

\[
\phi(X) = \beta_1 \text{Stab}(G_{\mathcal{P}}) + \beta_2 \text{GoalAcc}(G_{\mathcal{P}}) + \beta_3 \text{Robust}(G_{\mathcal{P}}).
\]

\subsection{Category-Theoretic Angle}
Represent praxicon as a functor \(F: \mathcal{C}_{\text{perc}} \to \mathcal{C}_{\text{motor}}\) mapping perceptual to motor categories, with semantic mediation. EAIMR measures natural transformations quantifying improbability of mappings \citep{jaynes2003}.

\section{Worked Examples}
\label{sec:examples}

\subsection{Neural Tissue}
High \(A(X)\) due to synaptic organization, praxicon-like engrams in parietal-premotor circuits yield high \(C(X)\). Simulations show low action entropy \citep{vannuscorps2023}. EAIMR detects mind-like agency.

\textbf{Entropy Curve}: \(H_{\text{obs}}(\mathcal{P}) \ll H(\mathcal{P}_0)\).

\subsection{Termite Mound}
Moderate \(A(X)\), high \(C(X)\) from thermoregulation. Collective praxicon via labor motifs \citep{turner2000}. Graph: action nodes (digging) linked to semantics (ventilation).

\textbf{Data}: \cite{camazine2001} shows coordinated construction.

\subsection{Forest Ecosystem}
High \(A(X)\) in trophic networks, praxicon-like nutrient routing \citep{levin2019}. High degeneracy, yet coherence over centuries elevates EAIMR.

\textbf{Diagram}: Lattice of species interactions.

\subsection{Negative Case: Crystal Growth}
Low \(C(X)\), high \(D(X)\), resulting in low EAIMR due to simple physical rules, no praxicon-like structure \citep{hazen2007}.

\section{Applications}
\label{sec:applications}

\subsection{Neuroscience}
EAIMR-praxicon detects apraxia deficits and mirror neuron activity \citep{peigneux2004}. Action simulation models show praxicon grounding \citep{vannuscorps2023}.

\subsection{AI and Robotics}
Praxicon-like dictionaries in robot learning reduce action entropy. EAIMR flags emergent agency in embodied AI tasks \citep{pastra2011}.

### Ecological and Collective Minds}
Termite mounds and forests exhibit distributed praxicons, with EAIMR detecting collective intelligence \citep{camazine2001, levin2019}.

### Cross-Cultural Cognition}
Rituals and dances encode praxicon-like action repositories, detectable via EAIMR \citep{barsalou2008}.

\section{Limits and Future Research}
\label{sec:limits}

EAIMR requires computable proxies for \(K(X)\). Praxicon models need validation across non-human systems. Future work includes integrating with RSVP field equations and testing in AI alignment scenarios.

\section{Ethical and Philosophical Implications}
\label{sec:ethics}

EAIMR-praxicon avoids anthropocentric bias, supporting interspecies communication and AI rights \citep{dennett1987}. It aligns with process philosophy, viewing minds as dynamic assemblies \citep{whitehead1929}. Ecological stewardship benefits from recognizing collective minds \citep{simondon2017}.

\section{Decision Rule}
\label{sec:decision}

Threshold \(\tau\) classifies systems with \(\text{EAIMR}(X) \geq \tau\) as mind-like, using proxies like compression ratios and resilience curves.

\section{Conclusion}
\label{sec:conclusion}

The EAIMR-praxicon framework provides a robust, information-theoretic approach to detecting mind-like agency, bridging cognitive science, ecology, and AI. It advances our understanding of intelligence without requiring human-like narrative capacities.

\appendix

\section{Derivations and Pseudocode}
\label{app:derivations}

\subsection{EAIMR Derivation}
\[
S(X) = I_0(X) + \lambda \phi_{\mathcal{P}}(X),
\]
\[
\text{EAIMR}(X) = g(S(X)), \quad g \uparrow.
\]

\subsection{Pseudocode for EAIMR}
\begin{verbatim}
function compute_eaimr(X, P0, P1, G_P):
    A = compute_assembly_index(X)
    R = estimate_reuse(X, motif_library)
    P = estimate_parallelism(X, concurrency_graph)
    D = estimate_degeneracy(X, functional_equivalents)
    C = compute_coherence(G_P, stability, goal_accuracy, robustness)
    return A / (R * P * D) * C
\end{verbatim}

\bibliographystyle{plainnat}
\bibliography{mind_at_work}

\end{document}